\chapter{Dokumentasi Praktikum}

\section{Hasil Praktikum}

\subsection{Inisiasi Komponen dan Fungsi Dasar}

\begin{figure}[H]
    \centering
    \includegraphics[width=0.6\textwidth]{1.png}
    \caption{Pembuatan Proyek Baru}
    \label{fig:1}
\end{figure}
Keterangan Gambar :
\begin{itemize}
    \item Membuat proyek Unreal Engine baru menggunakan template \textbf{Third Person} dengan penamaan yang sesuai format.
\end{itemize}

\begin{figure}[H]
    \centering
    \includegraphics[width=0.6\textwidth]{2.png}
    \caption{Menambahkan Blueprint Class Baru}
    \label{fig:2}
\end{figure}
Keterangan Gambar :
\begin{itemize}
    \item Membuat Blueprint Class baru dengan \textit{Parent Class} berupa \textbf{Actor Component}, kemudian diberi nama \textbf{BPC\_Health}.
\end{itemize}

\begin{figure}[H]
    \centering
    \includegraphics[width=0.6\textwidth]{3.png}
    \caption{Membuat Function ModifyHealth}
    \label{fig:3}
\end{figure}
Keterangan Gambar :
\begin{itemize}
    \item Menambahkan sebuah fungsi baru bernama \textbf{ModifyHealth} di dalam \textit{BPC\_Health} untuk menangani logika perubahan darah.
\end{itemize}

\begin{figure}[H]
    \centering
    \includegraphics[width=0.6\textwidth]{4.png}
    \caption{Menambahkan Variabel Health}
    \label{fig:4}
\end{figure}
Keterangan Gambar :
\begin{itemize}
    \item Membuat dua buah variabel bertipe \textit{Float} pada komponen, yaitu \textbf{MaxHealth} dan \textbf{CurrentHealth}.
\end{itemize}

\begin{figure}[H]
    \centering
    \includegraphics[width=0.6\textwidth]{5.png}
    \caption{Pembuatan Event Dispatcher}
    \label{fig:5}
\end{figure}
Keterangan Gambar :
\begin{itemize}
    \item Menambahkan \textit{Event Dispatcher} baru dengan nama \textbf{OnDeath} sebagai sinyal pemicu (trigger) saat karakter kehabisan darah.
\end{itemize}

\begin{figure}[H]
    \centering
    \includegraphics[width=0.6\textwidth]{6.png}
    \caption{Event Graph BPC\_Health}
    \label{fig:6}
\end{figure}
Keterangan Gambar :
\begin{itemize}
    \item Tampilan awal lembar kerja \textit{Event Graph} pada \textbf{BPC\_Health} sebelum diberikan node logika.
\end{itemize}

\begin{figure}[H]
    \centering
    \includegraphics[width=0.6\textwidth]{7.png}
    \caption{Logika Begin Play pada BPC\_Health}
    \label{fig:7}
\end{figure}
Keterangan Gambar :
\begin{itemize}
    \item Pada saat permainan dimulai (\textit{BeginPlay}), nilai \textbf{CurrentHealth} akan disetel (\textit{Set}) agar sama dengan nilai \textbf{MaxHealth} yang default-nya telah diatur menjadi 100.
\end{itemize}

\begin{figure}[H]
    \centering
    \includegraphics[width=0.6\textwidth]{8.png}
    \caption{Logika pada Function ModifyHealth}
    \label{fig:8}
\end{figure}
Keterangan Gambar :
\begin{itemize}
    \item Mendefinisikan alur pada fungsi \textbf{ModifyHealth} dengan menambahkan variabel input \textit{Health} (nilai 100) yang bertugas mengurangi darah karakter.
\end{itemize}

\subsection{Integrasi dengan Karakter dan Pembuatan Interface}

\begin{figure}[H]
    \centering
    \includegraphics[width=0.6\textwidth]{9.png}
    \caption{Tampilan Awal Event Graph Karakter Utama}
    \label{fig:9}
\end{figure}
Keterangan Gambar :
\begin{itemize}
    \item Membuka \textit{Event Graph} bawaan dari Blueprint karakter utama (\textbf{BP\_ThirdPersonCharacter} atau BP\_TPP).
\end{itemize}

\begin{figure}[H]
    \centering
    \includegraphics[width=0.6\textwidth]{10.png}
    \caption{Menambahkan Event OnDeath}
    \label{fig:10}
\end{figure}
Keterangan Gambar :
\begin{itemize}
    \item Memanggil node event \textbf{OnDeath} ke dalam lembar kerja karakter untuk bersiap mengeksekusi aksi saat kondisi mati tercapai.
\end{itemize}

\begin{figure}[H]
    \centering
    \includegraphics[width=0.6\textwidth]{11.png}
    \caption{Pembuatan Blueprint Interface}
    \label{fig:11}
\end{figure}
Keterangan Gambar :
\begin{itemize}
    \item Membuat kelas \textit{Blueprint Interface} baru di Content Drawer dan memberinya nama \textbf{BPI\_Overlap}.
\end{itemize}

\begin{figure}[H]
    \centering
    \includegraphics[width=0.6\textwidth]{12.png}
    \caption{Daftar Blueprint Proyek}
    \label{fig:12}
\end{figure}
Keterangan Gambar :
\begin{itemize}
    \item Tampilan seluruh aset Blueprint yang telah berhasil dibuat untuk keperluan sistem status karakter pada praktikum ini.
\end{itemize}

\begin{figure}[H]
    \centering
    \includegraphics[width=0.6\textwidth]{13.png}
    \caption{Pemasangan Komponen Health}
    \label{fig:13}
\end{figure}
Keterangan Gambar :
\begin{itemize}
    \item Memasukkan \textbf{BPC\_Health} ke dalam daftar komponen milik \textbf{BP\_TPP} agar sang karakter resmi memiliki modul sistem darah.
\end{itemize}

\begin{figure}[H]
    \centering
    \includegraphics[width=0.6\textwidth]{14.png}
    \caption{Logika Sederhana Event OnDeath}
    \label{fig:14}
\end{figure}
Keterangan Gambar :
\begin{itemize}
    \item Menyusun node instruksi ketika \textit{OnDeath} terpicu dan melakukan \textit{bind} pada \textbf{BP\_TPP}.
\end{itemize}

\begin{figure}[H]
    \centering
    \includegraphics[width=0.6\textwidth]{15.png}
    \caption{Fungsi TakeDamage pada Interface}
    \label{fig:15}
\end{figure}
Keterangan Gambar :
\begin{itemize}
    \item Di dalam \textbf{BPI\_Overlap}, ditambahkan fungsi \textbf{TakeDamage} dengan parameter input bernama \textbf{AmountDamage}.
\end{itemize}

\begin{figure}[H]
    \centering
    \includegraphics[width=0.6\textwidth]{16.png}
    \caption{Implementasi Interface pada Karakter}
    \label{fig:16}
\end{figure}
Keterangan Gambar :
\begin{itemize}
    \item Membuka menu \textit{Class Settings} pada karakter, lalu menambahkan \textbf{BPI\_Overlap} di bagian \textit{Implemented Interfaces}.
\end{itemize}

\begin{figure}[H]
    \centering
    \includegraphics[width=0.4\textwidth]{17.png}
    \caption{Verifikasi Implementasi Interface}
    \label{fig:17}
\end{figure}
Keterangan Gambar :
\begin{itemize}
    \item Indikator bahwa antarmuka berhasil terpasang, ditandai dengan munculnya fungsi \textbf{Take Damage} pada panel fungsi milik karakter.
\end{itemize}

\begin{figure}[H]
    \centering
    \includegraphics[width=0.6\textwidth]{18.png}
    \caption{Logika Implementasi Fungsi Interface}
    \label{fig:18}
\end{figure}
Keterangan Gambar :
\begin{itemize}
    \item Menerapkan logika saat \textit{Event TakeDamage} terpanggil. Fungsi \textbf{ModifyHealth} dijalankan, lalu nilai darah terbaru akan dicetak ke layar melalui \textit{Print String}.
\end{itemize}

\subsection{Pembuatan Aktor Pemicu Tabrakan (Collision)}

\begin{figure}[H]
    \centering
    \includegraphics[width=0.6\textwidth]{19.png}
    \caption{Pembuatan BP\_Collision}
    \label{fig:19}
\end{figure}
Keterangan Gambar :
\begin{itemize}
    \item Membuat aktor baru bernama \textbf{BP\_Collision}. Opsi \textbf{Hidden in Game} pada \textit{Box Collision}-nya dinonaktifkan agar garis panduannya terlihat saat dimainkan.
\end{itemize}

\begin{figure}[H]
    \centering
    \includegraphics[width=0.6\textwidth]{20.jpg}
    \caption{Logika Event Graph pada BP\_Collision}
    \label{fig:20}
\end{figure}
Keterangan Gambar :
\begin{itemize}
    \item Menyusun logika \textit{Eventgraph}. Saat box tersentuh, ia mengirimkan node \textbf{Take Damage (Message)} ke aktor yang menabraknya (\textit{Other Actor}).
\end{itemize}

\begin{figure}[H]
    \centering
    \includegraphics[width=0.6\textwidth]{21.png}
    \caption{Pengaturan Variabel Instance Editable}
    \label{fig:21}
\end{figure}
Keterangan Gambar :
\begin{itemize}
    \item Opsi \textbf{Instance Editable} pada variabel \textit{Damage} diaktifkan agar parameternya terbuka untuk modifikasi cepat di \textit{level}.
\end{itemize}

\begin{figure}[H]
    \centering
    \includegraphics[width=0.6\textwidth]{22.png}
    \caption{Mengubah Nilai Variabel Melalui Level}
    \label{fig:22}
\end{figure}
Keterangan Gambar :
\begin{itemize}
    \item Nilai kerusakan (\textit{Damage}) dari \textbf{BP\_Collision} kini dapat diubah-ubah secara spesifik melalui panel \textit{Details} di \textit{Level Editor}.
\end{itemize}

\clearpage
\section{Hasil Akhir}

\begin{figure}[H]
    \centering
    \includegraphics[width=0.6\textwidth]{23.png}
    \caption{Tampilan Karakter Sebelum Hit Collision}
    \label{fig:23}
\end{figure}
Keterangan Gambar :
\begin{itemize}
    \item Ketika simulasi di-play, karakter berada dalam kondisi normal.
    \item Terdapat tampilan log teks yang mencetak status darah pada layar (hasil \textit{BeginPlay}), yakni bernilai \textbf{100}.
\end{itemize}

\begin{figure}[H]
    \centering
    \includegraphics[width=0.6\textwidth]{24.png}
    \caption{Tampilan Karakter Setelah Hit Collision}
    \label{fig:24}
\end{figure}
Keterangan Gambar :
\begin{itemize}
    \item Karakter secara fisik masih normal setelah berinteraksi dengan kotak \textit{collision}.
    \item Terjadi penurunan jumlah darah setiap kali terjadinya sentuhan. Pada layar akan tercetak angka sisa darah terkini (misal \textbf{70} jika terjadi pengurangan 30 poin).
\end{itemize}

\begin{figure}[H]
    \centering
    \includegraphics[width=0.6\textwidth]{25.png}
    \caption{Kondisi Karakter Mati (OnDeath)}
    \label{fig:25}
\end{figure}
Keterangan Gambar :
\begin{itemize}
    \item Ketika karakter terus menerus menabrak kotak \textit{collision} hingga perhitungan fungsi \textit{Health} menyentuh angka \textbf{0}.
    \item Pemicu logika status karakter mati (\textbf{OnDeath}) berjalan, yang merespons dengan menampilkan cetakan string \textbf{"Tewas"} ke layar.
    \item Selain memunculkan teks, event \textbf{OnDeath} tersebut juga memicu perubahan ukuran secara visual, di mana skala karakter (Scale XYZ) akan membesar \textbf{10 kali lipat} dari ukuran aslinya.
\end{itemize}

\chapter{Daftar Pustaka}

\begin{enumerate}
    % Rujukan utama yang paling penting untuk Blueprint UE5
    \item \label{epic} Epic Games. \textit{Unreal Engine 5 Documentation}. [Online]. Tersedia: \url{https://docs.unrealengine.com/}.
    
    % Buku yang sangat spesifik membahas Node Blueprint UE5
    \item \label{romero} M. Romero. \textit{Blueprints Visual Scripting for Unreal Engine 5}, 3rd ed. Birmingham, UK: Packt Publishing, 2022.
    
    % Buku fundamental untuk pemahaman konsep game engine (seperti Event Tick, Begin Play)
    \item \label{gregory} J. Gregory. \textit{Game Engine Architecture}, 3rd ed. Boca Raton, FL: CRC Press, 2018.
    
    % Buku tambahan untuk logika Blueprint dasar
    \item \label{dorantes} M. Dorantes. \textit{Unreal Engine for Beginners: Learn to create games with Blueprint}. Independently published, 2020.
    
    % Fundamental grafis/rendering untuk engine (opsional tapi bagus untuk memperkuat laporan)
    \item \label{akenine} T. Akenine-Möller, E. Haines, dan N. Hoffman. \textit{Real-Time Rendering}, 4th ed. Boca Raton, FL: A K Peters/CRC Press, 2018.
\end{enumerate}